\scanpage{657}\mspage{45}
dans Hom. Alg. -- à cela près qu'il \deleted{\unsure{en}} prend
$D^-(K)$ au lieu de $D^b(k)$. Peut-être,
pour vérifier les axiomes, doit-on supposer
que $k$ est de dim. coh. finie.

Prenons $I=B_G$, où $G$ est un
groupe (p.ex. $G=\{\pm1\}$), \deleted{Alors} et
$u:I\to e$. Alors $M(I)$ est la
catégorie des complexes bornés de
$k[G]$-modules. On voit que
$u^W_!$ resp. $u^W_*$ existe \unsure{sss|ssi} pour tt
tel complexe borné, ou simplement pour
un $k[G]$-module sans plus $M$,
le complexe $L\unsure{\Gamma|\mathscr T}^{G}(M)$ resp. $R\Gamma_G(M)$
\deleted{$\unsure{L\Gamma^{G}(M)}$} appartient à $D^b(k)$, ou ce
qui revient au $=$, sss on a
\[
 \begin{aligned}
 H_i(G,M)&=0 &&\text{pour }i\gg0\\
 \text{\deleted{\unsure{resp.}}}\quad H^i(G,M)&=0 &&\text{pour }i\gg0,
 \end{aligned}
\]
\unsure{tout} $k[\unsure{\ell|G|k}]$-module $M$. Ce n'est
pratiquement jamais le cas ; \deleted{pour} quand
on prend p.ex. $k=\mathbb Z$ et $G$ groupe fini,
cela n'est le cas que si $G=e$. L'exemple
le plus économique \deleted{\unsure{est}} \inserted{\unsure{paraît|serait}} $G=\pm1$, $k=\mathbb F_2$.
% archive sheet number: 656

